%% kapa: exam=1 round=37 year=2026 period=2 form=A subject=appraisal_law subject_ko=감정평가관계법규 seq=25 number=25
%% source: Q-Net 감정평가사 기출문제 (cst00309) artlSeq=5258924 file=raw/1008/5258924/2026년 제37회 감정평가사 1차 2교시.pdf
%% conversion: kordoc 4.24.51 -> tools/gen1.py (auto); review=proofread-sonnet; flags=none
\begin{question}{25}
건축법령상 허가 대상 건축물이라 하더라도 건축신고를 하면 건축허가를 받은 것으로 보는 경우를 모두 고른 것은? (단, 주어진 조건 외에 다른 조건은 고려하지 않음)
\begin{framedtext}
ㄱ. 연면적이 160제곱미터이고 2층인 건축물의 대수선 \par
ㄴ. ｢국토의 계획 및 이용에 관한 법률｣에 따른 지구단위계획구역에서 연면적이 180제곱미터이고 2층인 건축물의 건축 \par
ㄷ. 소규모 건축물로서 연면적의 합계가 120제곱미터인 건축물의 건축 \par
ㄹ. 2층 건축물에서 바닥면적의 합계가 100제곱미터인 증축 \par
\end{framedtext}
\begin{choices}
  \item ㄱ
  \item ㄱ, ㄷ
  \item ㄴ, ㄹ
  \item ㄴ, ㄷ, ㄹ
  \item ㄱ, ㄴ, ㄷ, ㄹ
\end{choices}
\end{question}
\answer{1}
